\section{Matemáticas y Simulación}

\subsection{Patrón de Reducción Binaria (tipo Josefo)}

Preguntas guía — ¿cuándo usarlo?

\begin{itemize}
\item ¿El problema describe un proceso repetitivo/alternado (cada segundo elemento se "marca") sobre n elementos, con n muy grande (≥10\textasciicircum{}9)?
\item ¿Simular paso a paso sería O(n), demasiado lento, pero el patrón se repite de forma autosimilar al dividir por la mitad?
\item ¿Puedo replantear el problema de tamaño n como un subproblema de tamaño \textasciitilde{}n/2 aplicando la misma regla?
\end{itemize}

¿Para qué sirve? Resolver en O(log n) problemas de simulación con patrones alternados sobre secuencias enormes, sin simular cada paso. Es la misma idea del clásico Problema de Josefo: en vez de eliminar uno por uno, se "salta" de una ronda completa a la mitad del problema restante.

Complejidad: O(log n) por consulta (cada ronda aprox. divide n a la mitad).

Fundamento teórico: En cada "ronda" del proceso, la mitad de los elementos son descartados/activados según la paridad de su posición. Reformulando la posición y el tamaño del problema tras cada ronda (n' = n − procesados, i' = nueva posición relativa) se llega a una recurrencia que converge en O(log n) pasos en vez de O(n).

Ejercicio aplicado

Hay N personas en fila y un detector de metales defectuoso que se activa de forma alternada (la 1ª persona no lo activa, la 2ª sí, la 3ª no, ...; quien lo activa va al final de la fila). Dado que el coach está en la posición i, determinar en qué turno (entre quienes no activan el detector) pasará.

E/S: T ≤ 1000 casos, cada uno con solo dos enteros → readline por línea (patrón 1) es suficiente y más simple.

Código solución

\begin{lstlisting}[language=Python]
import sys
input = sys.stdin.buffer.readline
 
def safe_position(n, i):
# turno "seguro" (que no activa el patron) en que pasara la posicion i
    res = 0
    while i % 2 == 0:          # mientras i cae en la mitad "activada"
        res += i // 2         # cuenta los turnos seguros de esta ronda
        n -= i // 2             # reduce el problema a la mitad restante
        i = n                   # la posicion i se convierte en el nuevo n
    res += (i + 1) // 2        # suma los turnos seguros de la ronda final
    return res
 
def solve():
    t = int(input())
    out = []
    for _ in range(t):
        n, i = map(int, input().split())
        out.append(str(safe_position(n, i)))
    sys.stdout.write("\n".join(out))
 
solve()
\end{lstlisting}


\subsection{Fuerza Bruta / Simulación Directa}

Preguntas guía — ¿cuándo usarlo?

\begin{itemize}
\item ¿El tamaño de la entrada es pequeño (N ≤ 10\textasciicircum{}3 o similar) y no hay una fórmula cerrada evidente?
\item ¿Puedo calcular la respuesta iterando directamente sobre el espacio del problema sin que se dispare el tiempo?
\item ¿El problema es más de matemáticas/probabilidad que de estructuras de datos?
\end{itemize}

¿Para qué sirve? Reconocer cuándo NO hace falta un algoritmo sofisticado: con restricciones pequeñas, iterar directamente (simulando la definición del problema) es la solución más simple, rápida de programar y menos propensa a errores. Identificar esto ahorra tiempo valioso en el contrarreloj.

Complejidad: Depende del problema; en este caso O(n) con aritmética de punto flotante.

Fundamento teórico: Se calcula directamente la probabilidad de que NO haya coincidencia de cumpleaños entre k personas, multiplicando factores (1 − (k−1)/n) sucesivamente (extensión del problema del cumpleaños), deteniéndose apenas esa probabilidad cae por debajo de 0.5 (la probabilidad de coincidencia supera el 50\%).

Ejercicio aplicado

Dado N (número de días del año), hallar el mínimo número de personas necesarias en una sala para que la probabilidad de que dos compartan cumpleaños supere el 50\%.

E/S: Entrada de un solo entero → readline (patrón 1), no requiere nada especial.

Código solución

\begin{lstlisting}[language=Python]
import sys
input = sys.stdin.buffer.readline
 
n = int(input())
prob_no_coincidencia = 1.0
for k in range(1, n + 2):
    prob_no_coincidencia *= (n - (k - 1)) / n
    if prob_no_coincidencia < 0.5:
        print(k)
        break
\end{lstlisting}


\subsection{Simulación Lineal (Máximo de Prefijo)}

Preguntas guía — ¿cuándo usarlo?

\begin{itemize}
\item ¿La respuesta depende de un solo recorrido lineal del arreglo, calculando un máximo/mínimo/acumulado por posición?
\item ¿No hay dependencias entre posiciones más allá de una comparación directa fila por fila?
\item ¿El problema, aunque tenga una narrativa elaborada, se reduce matemáticamente a una fórmula simple tras identificar la variable clave?
\end{itemize}

¿Para qué sirve? Resolver directamente problemas donde, tras modelar bien la situación (aquí: el "cuello de botella" es la fila con menor espacio libre entre dos paredes), un solo recorrido O(n) calculando un máximo/mínimo basta. Es clave en competencia reconocer cuándo un problema de apariencia compleja se reduce a esto.

Complejidad: O(h) tiempo, O(h) espacio (o O(1) extra si se procesa en streaming).

Fundamento teórico: Las dos paredes se mueven a 1 m/s cada una. En la fila i, el espacio libre es w − (l[i] + r[i]) y se cierra a razón de 2 m/s (ambas paredes se acercan). El tiempo en que la fila i se cierra es (w − l[i] − r[i]) / 2. Las paredes chocan en la fila que se cierra primero: la de menor espacio libre inicial, equivalente al máximo de (l[i] + r[i]).

Ejercicio aplicado

Una sala de H filas y ancho W tiene una pared este y una pared oeste que se acercan a 1 m/s cada una desde que alguien entra. Cada fila i tiene un ancho ya ocupado por cada pared. Determinar en cuánto tiempo chocan las paredes.

E/S: H ≤ 10\textasciicircum{}5, dos arreglos de tamaño H → lectura total + tokenización (patrón 2).

Código solución

\begin{lstlisting}[language=Python]
import sys
 
def solve():
    data = sys.stdin.buffer.read().split()
    h = int(data[0]); w = int(data[1])
    left = list(map(int, data[2:2 + h]))
    right = list(map(int, data[2 + h:2 + 2 * h]))
    max_sum = max(left[i] + right[i] for i in range(h))
    ans = (w - max_sum) / 2
\end{lstlisting}


\begin{lstlisting}[language=Python]
if ans.is_integer():        # verifica parseo de float a int
    print(int(ans))
else:
    print(ans)
solve()
\end{lstlisting}


\subsection{+Ordenamiento de Puntos Críticos y Búsqueda de Huecos-----}

Preguntas guía — ¿cuándo usarlo?

\begin{itemize}
\item ¿El problema depende de encontrar el mayor (o menor) espacio libre entre eventos/puntos dispersos en una línea de tiempo o recta?
\item ¿Los puntos relevantes (incluyendo extremos del rango) se pueden ordenar y luego recorrer una sola vez?
\item ¿La respuesta se reduce a comparar diferencias entre elementos consecutivos ya ordenados?
\end{itemize}

¿Para qué sirve? Determinar si existe un hueco suficientemente grande entre puntos dispersos (tiempos, posiciones) en un intervalo, agregando los extremos del intervalo como puntos artificiales y ordenando todo junto. Evita simular minuto a minuto: basta con mirar diferencias consecutivas tras ordenar.

Complejidad: O(n log n) por el ordenamiento; O(n) el resto.

Fundamento teórico: Cualquier hueco maximal sin eventos queda delimitado por dos puntos consecutivos en el arreglo ordenado (incluyendo los extremos del intervalo total como puntos ficticios). Por lo tanto, el hueco más grande siempre es la mayor diferencia entre elementos consecutivos de la lista ordenada; no hace falta considerar ninguna otra combinación de puntos.

Ejercicio aplicado

Un vuelo dura D minutos y sirve M comidas en instantes dados; hace falta dormir T minutos seguidos sin interrupción. Determinar si es posible dormir T minutos consecutivos sin perderse ninguna comida.

E/S: M ≤ 1000 → entrada pequeña, readline por línea (patrón 1) es suficiente.

Código solución

\begin{lstlisting}[language=Python]
import sys
input = sys.stdin.buffer.readline
 
def solve():
    T, D, M = map(int, input().split())
    times = [0, D]                      # extremos del vuelo como puntos ficticios
    for _ in range(M):
        y = int(input())
        times.append(y)
    times.sort()
 
    ans = 'N'
    for i in range(1, len(times)):
        if times[i] - times[i - 1] >= T:   # hueco suficiente para dormir T minutos
            ans = 'Y'
            break
    print(ans)
 
solve()
\end{lstlisting}


\subsection{+Teoría de Números: Regla de Divisibilidad Alternada----}

Preguntas guía — ¿cuándo usarlo?

\begin{itemize}
\item ¿El problema pide verificar o forzar divisibilidad por (base ± 1) de un número representado por dígitos?
\item ¿Se puede aprovechar que base ≡ −1 (mod base+1), lo que convierte la suma ponderada de dígitos en una suma ALTERNADA módulo (base+1)?
\item ¿Se pide además construir/ajustar un dígito específico para lograr esa divisibilidad?
\end{itemize}

¿Para qué sirve? Generalizar la clásica regla de divisibilidad por 11 en base 10 (suma alterna de dígitos) a cualquier base B, para verificar o forzar divisibilidad por B+1 sin operar con el número completo (que puede tener hasta 2·10\textasciicircum{}5 dígitos).

Complejidad: O(L), con L = cantidad de dígitos.

Fundamento teórico: Como B ≡ −1 (mod B+1), cada potencia B\textasciicircum{}i ≡ (−1)\textasciicircum{}i (mod B+1). Por lo tanto, el valor del número módulo B+1 es exactamente la suma alternada de sus dígitos (sumando los de posición par, restando los de posición impar). Para forzar que esa suma sea 0 módulo B+1 cambiando un solo dígito D\_i por un valor MENOR, se despeja cuánto debe valer ese dígito para anular el residuo actual, y se verifica que el nuevo valor sea válido (entre 0 y el dígito original, sin aumentarlo).

Ejercicio aplicado

Dado un número N de L dígitos en base B, encontrar un dígito que se pueda reducir (nunca aumentar) para que el número resultante sea múltiplo de B+1, priorizando el menor valor final posible, o determinar que no existe solución.

E/S: L ≤ 2·10\textasciicircum{}5 → aunque el código usa input() estándar por simplicidad (ya venía así y funciona), para esta escala también sería válido leer con sys.stdin.buffer.read().split() (patrón 2) si se buscara el máximo rendimiento.

Código solución

\begin{lstlisting}[language=Python]
# Lee la base B y la cantidad de digitos L
basedig = list(map(int, input().split()))
nums = list(map(int, input().split()))   # digitos D_1..D_L, de mas a menos significativo
 
base = basedig[0]
div = basedig[0] + 1                     # queremos que N sea multiplo de base+1
dig = basedig[1]                         # L = cantidad de digitos
 
# suma alternada de digitos == N mod (base+1), porque base = -1 mod (base+1)
n = 0
for i in range(len(nums)):
    if (dig - 1 - i) % 2 == 0:
        n += nums[i]
    else:
        n -= nums[i]
n %= div
 
if n == 0:
    print("0 0")                         # ya es multiplo, no hace falta cambiar nada
    exit()
 
# buscar el primer digito (de izquierda a derecha) que se pueda REDUCIR
# para anular el residuo n, segun la paridad de su posicion
res = 0
for i in range(len(nums)):
    if (dig - 1 - i) % 2 == 0:
        res = n                          # esta posicion suma: hay que restarle 'n'
    else:
        res = div - n                    # esta posicion resta: se compensa distinto
    if 1 <= res <= nums[i]:              # solo vale si es una reduccion real (no negativa)
        print(f"{i + 1} {nums[i] - res}")
        exit()
 
print("-1 -1")                           # ningun digito puede arreglarlo
\end{lstlisting}


\subsection{+Simulación de Eventos con Colas------}

Preguntas guía — ¿cuándo usarlo?

\begin{itemize}
\item ¿Hay un recurso compartido con estado (dirección, disponibilidad) que cambia según el orden de llegada de eventos con timestamps?
\item ¿Los eventos ya vienen ordenados (o se pueden ordenar) y conviene procesarlos con colas separadas por tipo?
\item ¿El estado siguiente depende de comparar el próximo evento contra un tiempo límite (fin de uso actual del recurso)?
\end{itemize}

¿Para qué sirve? Simular un sistema con estado (aquí: una dirección de movimiento y un instante de fin de uso) que evoluciona a medida que llegan eventos de distintos tipos, usando una cola por tipo de evento en vez de simular minuto a minuto.

Complejidad: O(n) — cada persona se encola y desencola una sola vez.

Fundamento teórico: Se mantienen dos colas (una por dirección deseada) y un estado (dirección activa, instante en que el recurso queda libre). En cada paso se mira el evento más próximo entre ambas colas: si coincide con la dirección activa, extiende el fin de uso; si es la dirección contraria, esa persona debe esperar (se cuenta pero no se procesa aún); si ya no hay eventos antes de que el recurso quede libre, se decide si invertir la dirección (si había gente esperando) o terminar ese ciclo de uso.

Ejercicio aplicado

N personas llegan a una escalera doble (10 segundos por cruce) en tiempos distintos, cada una queriendo ir en una de dos direcciones. Determinar el instante en que la última persona termina de usarla.

E/S: n ≤ 10\textasciicircum{}4 → lectura total + tokenización (patrón 2); se usan dos deques (collections.deque) en vez de queue, ya que en Python deque tiene O(1) real para popleft().

Código solución

\begin{lstlisting}[language=Python]
import sys
from collections import deque
 
def solve():
    data = sys.stdin.buffer.read().split()
    idx = 0
    n = int(data[idx]); idx += 1
    left = deque()    # personas que quieren ir en direccion 0
    right = deque()   # personas que quieren ir en direccion 1
    for _ in range(n):
        t = int(data[idx]); d = int(data[idx + 1]); idx += 2
        (left if d == 0 else right).append(t)
 
    end = 0          # instante en que el recurso queda libre
    ans = 0
    waiting = 0       # cuantos esperan en direccion contraria a la activa
    direction = -1    # -1 = detenida
 
    while left or right or direction != -1:
        # mira el proximo evento disponible entre ambas colas
        next_time, who = None, -1
        if left and (next_time is None or left[0] < next_time):
            next_time, who = left[0], 0
        if right and (next_time is None or right[0] < next_time):
            next_time, who = right[0], 1
 
        if direction == -1:
            # arranca detenida: la toma quien llegue primero
            direction = who
            end = next_time + 10
            (left if who == 0 else right).popleft()
        elif who != -1 and next_time < end:
            if who == direction:
                end = max(end, next_time + 10)   # se sube mientras sigue en marcha
            else:
                waiting += 1                     # direccion contraria: debe esperar
            (left if who == 0 else right).popleft()
        else:
            # no hay mas eventos antes de que el recurso quede libre
            if waiting > 0:
                direction = 1 - direction   # invierte sentido para atender a los que esperan
                end += 10
                waiting = 0
            else:
                ans = end
                direction = -1               # se detiene
 
    print(ans)
 
solve()
\end{lstlisting}
